\seccionreporte{Resultados}{sec:resultados}

\propositoseccion{%
  Presente métricas del modelo, comportamiento del servicio y ejemplos interpretables para la app móvil.%
}

\subsection{Resultados del modelo}

Como comentamos antes, HDBSCAN nos arrojó resultados bastante buenos con el corpus de 50 proyectos. El Silhouette de 0.69 nos dice que el modelo sí logró entender de qué trata cada documento y agruparlos lógicamente sin revolver turismo con fintech, por ejemplo. 

\subsection{Ejemplos de inferencia}

Aquí les dejo un ejemplo de lo que nos regresa el endpoint cuando le pegamos con un archivo real. Fíjense cómo el LLM nos da un veredicto súper estructurado.

\begin{lstlisting}[style=jsonstyle]
{
  "indice_innovacion": 85,
  "metricas_calidad": {
    "rigor_academico": 88,
    "relevancia_tecnica": 82
  },
  "riesgo_colision": {
    "porcentaje": 12.5,
    "nivel": "Bajo"
  },
  "recomendaciones_ia": [
    {
      "titulo": "Enfoque original",
      "descripcion": "El uso de IoT para este sector específico no tiene precedentes en el corpus histórico de la universidad."
    }
  ]
}
\end{lstlisting}

\subsection{Desempeño del servicio}

Hicimos algunas pruebas de carga ligeras para ver qué tanto aguantaba el servidor de Contabo (que tiene 8GB de RAM), y estos son los números:

\begin{tabular}{@{}lr@{}}
  \toprule
  Métrica & Valor \\
  \midrule
  Latencia media (Análisis Rápido HDBSCAN) & $\sim 1.5$ segundos \\
  Latencia media (Análisis Profundo Ollama/Phi-3) & $\sim 45$ segundos \\
  Tamaño de los modelos .joblib (UMAP + HDBSCAN) & $< 150$ KB (ligerísimo) \\
  Consumo de RAM de Ollama en inferencia & $\sim 4.2$ GB \\
  \bottomrule
\end{tabular}

\subsection{Visualizaciones}

\figuraopcional{figures/mapa_semantico_2d.png}{0.8\textwidth}{%
  El mapa semántico muestra claramente cómo los proyectos de una misma categoría se amontonan, lo que confirma que el modelo de embeddings funciona bien.%
}{Visualización 2D de clústeres UMAP}

\subsection{Limitaciones conocidas}

Obviamente no todo es perfecto, detectamos algunas cosillas que tenemos que tener en cuenta:

\begin{itemize}
  \item \textbf{RAM del Servidor:} Como Ollama se come casi 4GB de RAM cada vez que hace una inferencia, si dos alumnos le dan click al mismo tiempo, el contenedor de Docker hace cuello de botella.
  \item \textbf{Alucinaciones del LLM:} Aunque le pasamos los datos duros de colisión a Phi-3 para que no invente cosas, a veces en las "recomendaciones" tira tecnologías que no vienen al caso.
  \item \textbf{Textos en imágenes:} Si el alumno sube un PDF que es un escaneo (imágenes en vez de texto), nuestro parser falla porque no le metimos OCR para no alentar el proceso.
\end{itemize}
